%!TEX root = ../thesis.tex

\cleardoublepage
\phantomsection
\addcontentsline{toc}{chapter}{Part I: Structure-Aligned Learning for Visual Understanding}
\markboth{Structure-Aligned Learning for Visual Understanding}{Structure-Aligned Learning for Visual Understanding}
\thispagestyle{plain}

\begin{tikzpicture}[remember picture,overlay]
\node[align=center,text width=0.8\textwidth,inner sep=0pt] at (current page.center) {%
{\Huge\bfseries Part I\\Structure-Aligned Learning for\\Visual Understanding\par}
\vspace{1.5cm}
The next three chapters instantiate structure-aligned learning in panoptic scene graph generation (PSG). Chapter~3 makes the head--tail organization of predicates operational through complementary relation views. Chapter~4 aligns language-derived relational knowledge with visual evidence. Chapter~5 changes the pair representation and decoding procedure to move beyond a fixed predicate vocabulary. Together, these chapters develop the visual-understanding part of the argument in Section~\ref{sec:unifying_argument}.
};
\end{tikzpicture}
\null

\clearpage
